\documentclass[10pt,letterpaper]{article}
\usepackage[T1]{fontenc}
\usepackage[utf8]{inputenc}
\usepackage{lmodern}
\usepackage[margin=0.76in]{geometry}
\usepackage[table]{xcolor}
\usepackage{array,tabularx,longtable,booktabs}
\usepackage{enumitem}
\usepackage{amssymb,amsmath}
\usepackage{hyperref}
\usepackage{fancyhdr,lastpage}
\usepackage{microtype}
\usepackage{parskip}
\definecolor{cdblue}{HTML}{203957}
\definecolor{cdgray}{HTML}{586677}
\hypersetup{colorlinks=true,linkcolor=cdblue,urlcolor=cdblue,pdftitle={Chapter 03: Continuous Integration}}
\pagestyle{fancy}
\fancyhf{}
\lhead{\small\color{cdgray}Continuous Delivery Engineering}
\rhead{\small\color{cdgray}Study guide}
\cfoot{\thepage\ / \pageref{LastPage}}
\setlength{\headheight}{14pt}
\setlength{\emergencystretch}{2em}
\setlist[itemize]{leftmargin=*,itemsep=3pt,topsep=4pt}
\setlist[enumerate]{leftmargin=*,itemsep=3pt,topsep=4pt}
\begin{document}
\begin{center}
{\LARGE\bfseries\color{cdblue}Chapter 03: Continuous Integration}\\[7pt]
{\large I: Foundations}\\[6pt]
{\small Reference edition: 2026-10-30 $\bullet$ Original educational and operational material}
\end{center}
\noindent\color{cdblue}\rule{\linewidth}{1.2pt}\color{black}
\vspace{0.4em}
\section{Chapter scope and provenance}
Primary reading: Humble and Farley, Continuous Delivery (2010), Chapter 3, "Continuous Integration"; Part I: Foundations. This guide is an original synthesis and not a reproduction of the chapter. The book provides the foundational delivery principles; current AWS, Terraform and GitHub Actions references in the applied sections are supplementary.


\section{Learning outcomes}
\begin{enumerate}
\item Specify CI entry criteria and commit-stage ownership.
\item Contrast integration frequency with branch duration.
\item Establish broken-build repair and revert limits.
\end{enumerate}

\section{Core conceptual model}
Focus areas: trunk-based integration; fast commit tests; broken build response; distributed development.

\begin{itemize}
\item Input: a version-controlled change, explicit requirements, and a reproducible execution context.
\item Mechanism: controlled feedback, automation, evidence capture, and a stop condition where risk exceeds tolerance.
\item Output: a verifiable release candidate or a diagnosed failure with an accountable owner.
\end{itemize}

\section{Reading and analysis protocol}
\begin{enumerate}
\item Read the chapter and record five propositions, distinguishing an author claim from an example or assumption.
\item Identify two failure modes, one recovery technique, and one measurable acceptance condition.
\item Trace each technique to the stages and controls in the deployment-pipeline-architecture reference.
\item State how the proposed practice changes for a regulated system with segregated duties.
\end{enumerate}

\section{Detailed technical study notes}
\subsection{Mainline integration}
Short-lived changes reduce the graph of concurrent work that must be reconciled. Merge cadence affects the size of surprise integration failures.

\subsection{Commit stage contract}
A change entering mainline triggers compilation, deterministic tests and policy checks before downstream promotion.

\subsection{Broken-build ownership}
Treat failure as a team signal requiring containment. Record whether to fix, revert or temporarily halt integration.

\subsection{Fast tests and asynchronous checks}
Prioritize checks that return actionable feedback quickly; slow nonfunctional tests may follow without hiding commit-stage failures.

\subsection{CI system reproducibility}
The same tools and inputs must run locally and in CI where feasible; avoid network downloads that silently change dependency versions.

\subsection{Distributed collaboration}
Ensure developers in different regions see one shared source of truth and an auditable outcome for each integrated change.

\section{Domain-specific challenge questions}
\begin{enumerate}
\item For mainline integration, identify the precondition, expected evidence, adversarial failure case and measurable corrective action.
\item For commit stage contract, identify the precondition, expected evidence, adversarial failure case and measurable corrective action.
\item For broken-build ownership, identify the precondition, expected evidence, adversarial failure case and measurable corrective action.
\item For fast tests and asynchronous checks, identify the precondition, expected evidence, adversarial failure case and measurable corrective action.
\end{enumerate}
\section{Guided exercise}
Create a 5-minute commit pipeline with unit tests and a failing-test demonstration.

\subsection{Required evidence}
\begin{itemize}
\item A version-controlled design or executable artifact.
\item A failure demonstration, negative test or counterexample.
\item An improvement statement with baseline, target, owner and verification method.
\end{itemize}


\section{Engineering scenario}
A team uses multi-week branches and green feature-branch builds but continually breaks mainline.

\begin{itemize}
\item Describe the failure propagation path and the earliest feasible detection point.
\item Recommend a minimal remediation and its rollback or containment plan.
\item Define a measure that would demonstrate risk reduction.
\end{itemize}

\section{Knowledge checks}
\begin{enumerate}
\item Which condition is the gate, and which action is the response when it fails?
\item What evidence distinguishes successful automation from an assumed successful release?
\item Where might environments, data, dependencies or people invalidate the technique?
\item What feedback could be moved earlier without losing necessary confidence?
\item Which assumptions are book principles, and which are modern supplementary implementation choices?
\end{enumerate}

\section{Completion rubric}
Score each dimension from 0 (missing) to 4 (demonstrated and independently verifiable): conceptual accuracy; reproducibility; negative testing; security and safety; documentation and traceability. Pass threshold: 16/20 with no zero on reproducibility or safety.


\section{Prerequisites and onward route}
Git, basic tests, and a build runner. Next: Chapters 4 onward.



\end{document}
